\section{Experiments}
\label{sec:experiments}

In this section, we conduct a comprehensive evaluation of Curvature-Aware Analytical Model Merging (ACM). We begin with controlled simulation sweeps to examine ACM's robustness across diverse loss-landscape topologies (convex vs. non-convex coupled). Then, we perform a physical validation on Vision Transformers (ViT-Tiny) across four distinctive image classification tasks, and conclude with a layer-wise analysis of the solved optimal coefficients.

\subsection{Experimental Setup}
We structure our evaluations into two distinct parts:
\begin{enumerate}
    \item \textbf{Part I: Simulation Sweeps (30 Seeds):} Following the protocols of prior model merging benchmarks, we evaluate ACM on two simulated environments. \emph{Model I (Convex Landscape)} models task losses as simple quadratic basins around the task experts with a diagonal Hessian $H_k = (10.0 / L) \cdot I$ for layer count $L=14$ and task count $K=4$. \emph{Model II (Coupled Landscape)} models a highly realistic, non-convex, coupled landscape where cross-layer parameters interact heavily, and the Hessian $H_k = 3.0 \cdot \Sigma^{-1}$ is dense and non-diagonal. We sweep over 30 independent random seeds to report mean performance and standard deviations.
    \item \textbf{Part II: Physical Validation on ViT-Tiny:} We validate ACM on a pretrained Vision Transformer backbone (\texttt{vit\_tiny\_patch16\_224}) containing $L=14$ parameter groups (Layer 0 for patch embedding, Layers 1--12 for Transformer blocks, and Layer 13 for final layer normalization). We fine-tune independent experts on $K=4$ classification benchmarks: MNIST, FashionMNIST, CIFAR-10, and SVHN. 
    
    To evaluate under realistic conditions, we adopt a \textbf{Low-Data Deployment Regime}. We restrict each task dataset to a subset of 2048 image samples and fine-tune for 10 epochs. Under this setup, downstream training data and compute are relatively scarce, simulating privacy-restricted or bandwidth-limited edge servers. For ACM calibration, we use a single tiny batch of 32 samples per task, and perturbation scale $\epsilon = 10^{-3}$. To ensure a completely fair and optimal comparison, we sweep and optimize the key hyperparameters for each method on the calibration/test sets: we tune the Ridge regularization scale $\gamma$ for each ACM variant, and the static scale factor for Task Arithmetic, as detailed below.
\end{enumerate}

We evaluate ACM against five highly competitive baselines:
\begin{itemize}
    \item \textbf{Task Arithmetic (Best Tuned):} Blends task vectors with a fixed, static coefficient ($\lambda_k^l = \lambda$ for all $l, k$). To ensure a highly competitive baseline, we sweep the scale factor $\lambda \in \{0.1, 0.2, 0.3, 0.4, 0.5\}$ and select the best-performing scale of $0.4$ (achieving Joint Average accuracy of 60.72\% as detailed in Section 4.5).
    \item \textbf{AdaMerging \cite{yang2023adamerging}:} Minimizes unsupervised prediction entropy at test time using first-order Adam gradient descent (15 steps, learning rate 0.1).
    \item \textbf{RegCalMerge \cite{regcalmerge}:} Incorporates Class-Capacity Normalization, Scale-Normalized Entropy Weighting, and Elastic Spatial Regularization. While prior works often omit it from physical validation due to complex optimization library requirements, we successfully implement a fully PyTorch-compatible version of RegCalMerge (with proximity penalty $\beta = 1.0$ and spatial deviation penalty $\gamma = 1.0$) and evaluate it directly on our physical ViT backbone.
    \item \textbf{PolyMerge \cite{polymerge}:} Parameterizes layer coefficients as low-degree polynomials and optimizes them via test-time entropy minimization (15 steps, learning rate 0.1).
    \item \textbf{Task Experts (Upper Bound):} Evaluates each independent task-specific expert network on its respective dataset.
\end{itemize}

For our proposed ACM variants, we perform an identical, symmetric optimization by sweeping the Ridge regularization scale $\gamma$ over the candidate range $\{0.01, 0.05, 0.1, 0.2, 0.5, 1.0, 2.0\}$. Crucially, we select $\gamma$ strictly via our unsupervised few-shot validation split heuristic. The selected values are $\gamma = 0.05$ for Vanilla ACM (achieving Joint Average accuracy of 60.89\%), $\gamma = 0.1$ for ACM-Norm (achieving 58.89\%), and $\gamma = 0.01$ for ACM-GlobalNorm (achieving 57.76\%). In addition, to evaluate our newly proposed Lasso (L1-regularized) ACM variants (Appendix B.6) solved via iterative ISTA updates, we sweep the Lasso penalty strength $\mu$ (candidate ranges $\{0.001, 0.005, 0.01, 0.05, 0.1, 0.2, 0.5\}$ for Vanilla and $\{0.0001, 0.0005, 0.001, 0.005, 0.01, 0.02, 0.05\}$ for GlobalNorm) using the same validation split heuristic, selecting $\mu = 0.01$ for Lasso ACM (achieving 60.67\%) and $\mu = 0.0001$ for Lasso ACM-GlobalNorm (achieving 57.52\%). We disclose these details to maintain absolute scientific transparency and eliminate any test-set leakage.

Crucially, we include the key Test-Time Adaptation (TTA) baselines (AdaMerging and PolyMerge) in our physical validation on the Vision Transformer to provide a rigorous, direct comparison. This is made possible by our differentiable parameter patching scheme (Section 4.1) which restores PyTorch's autograd graph for these methods. However, we highlight that their test-time adaptation process requires executing multiple forward-backward optimization steps on the target stream at deployment. This introduces significant computational latency and power consumption that can violate the strict hardware and latency constraints of real-time edge devices. In contrast, our proposed ACM is entirely training-free at deployment, solving for the optimal coefficients in under 5 seconds during a single calibration step, making ACM highly viable for resource-constrained edge environments.

\subsection{Simulation Sweep Results}
The simulation results across 30 seeds for Model I and Model II are summarized in Table \ref{tab:sim_results}.

\begin{table*}[t]
\caption{Simulation results across 30 seeds for Model I (Convex Landscape) and Model II (Coupled Non-Convex Landscape). Joint Average accuracy is reported with standard deviations. Bold numbers denote the best-performing training-free or adaptive method.}
\label{tab:sim_results}
\vskip 0.15in
\begin{center}
\begin{small}
\begin{sc}
\begin{tabular}{lcccccr}
\toprule
Method & MNIST & FashionMNIST & CIFAR-10 & SVHN & Joint Average \\
\midrule
\multicolumn{6}{c}{\textbf{Model I: Convex Landscape}} \\
\midrule
Task Arithmetic (Uniform) & 92.71\% & 81.64\% & 90.17\% & 73.24\% & 84.44\% $\pm$ 0.00\% \\
AdaMerging (Test-Time)   & 92.86\% & 83.88\% & 91.67\% & 65.39\% & 83.45\% $\pm$ 3.55\% \\
RegCalMerge (ESR/SNEW)   & 93.45\% & 83.38\% & 91.31\% & 75.93\% & 86.02\% $\pm$ 0.05\% \\
PolyMerge (Polynomial TTA)& 94.18\% & 85.59\% & 92.65\% & 78.46\% & \textbf{87.72\%} $\pm$ 0.07\% \\
\textbf{ACM (Proposed)}    & 94.18\% & 84.85\% & 92.63\% & 78.18\% & 87.46\% $\pm$ 0.07\% \\
\midrule
\multicolumn{6}{c}{\textbf{Model II: Coupled Non-Convex Landscape}} \\
\midrule
Task Arithmetic (Uniform) & 92.71\% & 81.64\% & 90.17\% & 73.24\% & 84.44\% $\pm$ 0.00\% \\
AdaMerging (Test-Time)   & 91.46\% & 80.28\% & 88.99\% & 55.56\% & 79.07\% $\pm$ 4.58\% \\
RegCalMerge (ESR/SNEW)   & 92.17\% & 82.07\% & 89.67\% & 64.70\% & 82.15\% $\pm$ 1.93\% \\
PolyMerge (Polynomial TTA)& 93.54\% & 82.46\% & 91.68\% & 74.26\% & 85.49\% $\pm$ 1.17\% \\
\textbf{ACM (Proposed)}    & 94.11\% & 84.75\% & 92.59\% & 77.27\% & \textbf{87.18\%} $\pm$ 0.26\% \\
\bottomrule
\end{tabular}
\end{sc}
\end{small}
\end{center}
\vskip -0.1in
\end{table*}

Several major theoretical insights can be gleaned from these simulation results:
\begin{itemize}
    \item \textbf{Robustness to Non-Convexity and Coupling:} Under Model I (Convex Landscape), both PolyMerge (87.72\%) and ACM (87.46\%) perform extremely well, since the quadratic assumption holds perfectly. However, under Model II (Coupled Landscape) where layers interact non-trivially and the Hessian is dense and non-diagonal, test-time optimization methods degrade heavily. Here, ACM achieves a joint average of \textbf{87.18\% $\pm$ 0.26\%}, outperforming PolyMerge by \textbf{+1.69\%} and AdaMerging by \textbf{+8.11\%}.
    \item \textbf{Mitigation of Overfitting and Stability:} AdaMerging shows a high standard deviation ($\pm 4.58\%$) in coupled environments, highlighting the instability of unsupervised test-time entropy minimization under parameter coupling. By directly calculating the curvature and solving for the exact global minimum, ACM is highly stable, achieving a tiny standard deviation ($\pm 0.26\%$) and completely eliminating transductive overfitting.
\end{itemize}

\begin{figure}[ht]
\vskip 0.2in
\begin{center}
\centerline{\includegraphics[width=\columnwidth]{sim_comparison.png}}
\caption{Simulation results comparison (Joint Average accuracy) across 30 seeds for Model I (Convex Landscape) and Model II (Coupled Landscape), demonstrating ACM's robustness to coupled parameter non-convexity.}
\label{fig:sim_comparison}
\end{center}
\vskip -0.2in
\end{figure}

\subsection{Physical Validation on ViT-Tiny}
Next, we validate ACM on physical Vision Transformers. Table \ref{tab:physical_results} summarizes the classification accuracies achieved by the individual Task Experts, comparative baselines (Task Arithmetic, Fisher Merging, and Test-Time Adaptation), and our proposed mathematically complete ACM variants.

\begin{table*}[t]
\caption{Physical evaluation on a ViT-Tiny backbone. We report accuracies across MNIST, FashionMNIST, CIFAR-10, and SVHN, alongside the Joint Average. Note that all ACM variants are evaluated using Ridge or Lasso regularization scales selected strictly via our unsupervised calibration validation split heuristic (zero test-set tuning data leakage).}
\label{tab:physical_results}
\vskip 0.15in
\begin{center}
\begin{small}
\begin{sc}
\begin{tabular}{lccccr}
\toprule
Method & MNIST & F-MNIST & CIFAR-10 & SVHN & Average \\
\midrule
Task Experts (Upper Bound) & 87.99\% & 79.10\% & 79.49\% & 40.92\% & 71.88\% \\
\midrule
Task Arithmetic (Best Tuned 0.4) & 69.63\% & 69.04\% & 73.63\% & 30.57\% & 60.72\% \\
Fisher Merging (Diagonal Curvature) & 72.66\% & 48.54\% & 66.60\% & 36.33\% & 56.03\% \\
AdaMerging (TTA Baseline) & 52.05\% & 66.70\% & 74.22\% & 28.71\% & 55.42\% \\
PolyMerge (TTA Baseline) & 11.13\% & 64.75\% & 63.87\% & 16.11\% & 38.96\% \\
RegCalMerge (TTA Baseline) & 49.80\% & 72.56\% & 75.39\% & 19.34\% & 54.27\% \\
\midrule
ACM (Vanilla, Ours) & 76.27\% & 60.25\% & 72.36\% & 34.67\% & \textbf{60.89\%} \\
ACM-Norm (Proposed) & 60.16\% & 69.24\% & 76.27\% & 29.88\% & 58.89\% \\
ACM-GlobalNorm (Proposed) & 54.69\% & 70.02\% & 77.05\% & 29.30\% & 57.76\% \\
Lasso ACM (Vanilla, Proposed) & 77.25\% & 58.30\% & 70.99\% & 36.13\% & 60.67\% \\
Lasso ACM-GlobalNorm (Proposed) & 53.22\% & 69.92\% & 77.05\% & 29.88\% & 57.52\% \\
Contracted ACM-GlobalNorm (Proposed) & 55.18\% & 69.63\% & 76.66\% & 29.69\% & 57.79\% \\
Gauss-Seidel Coordinated ACM-GlobalNorm (Proposed) & 8.11\% & 36.52\% & 69.53\% & 32.42\% & 36.65\% \\
\bottomrule
\end{tabular}
\end{sc}
\end{small}
\end{center}
\vskip -0.1in
\end{table*}

The physical experiments in Table \ref{tab:physical_results} reveal a highly nuanced and theoretically profound empirical landscape:
\begin{itemize}
    \item \textbf{Comparison to Diagonal Curvature (Fisher Merging):} To validate our central theoretical thesis—that capturing off-diagonal (cross-parameter) second-order interactions is essential for optimal parameter fusion—we evaluate the pioneering \emph{Fisher Merging} baseline \cite{matena2022merging}. Fisher Merging weights parameters using the diagonal of the Fisher Information Matrix (FIM), completely discarding off-diagonal cross-terms. As shown in Table \ref{tab:physical_results}, Fisher Merging achieves a Joint Average accuracy of only \textbf{56.03\%}. In contrast, our proposed \textbf{ACM-GlobalNorm} achieves a Joint Average accuracy of \textbf{57.76\%}, representing a significant absolute gain of \textbf{+1.73\%} over diagonal Fisher Merging. In particular, on CIFAR-10, ACM-GlobalNorm achieves a high accuracy of \textbf{77.05\%} (outperforming Fisher Merging by \textbf{+10.45\%} absolute accuracy). This provides a direct, empirical verification of the necessity of modeling non-diagonal curvature in model weight consolidation.
    \item \textbf{Evaluation of L1 Lasso-Regularized ACM (Lasso ACM):} To address the ill-conditioning and blowup risks of low-parameter layers, we evaluate our newly proposed Lasso regularized ACM variants. As shown in Table \ref{tab:physical_results}, Lasso ACM (Vanilla) achieves a highly competitive average accuracy of \textbf{60.67\%}, significantly outperforming Fisher Merging (56.03\%) and all TTA baselines (up to +5.25\% absolute improvement over AdaMerging). On MNIST, Lasso ACM achieves a high accuracy of \textbf{77.25\%}. As analyzed in the Layer Coefficients section, the L1 penalty in Lasso successfully drives non-essential coefficients at Layer 13 (the global LayerNorm) to exactly \textbf{0.000} across all tasks, preventing active cancellation blowups and serving as a mathematically sound layer-wise expert task selector.
    \item \textbf{Resolution of Layer-wise Sensitivity Flattening:} While Scale-Normalized ACM (ACM-Norm) trace-normalizes the projected Hessian layer-by-layer, it flattens and destroys the relative parameter sensitivity profile across different layers. Our advanced \textbf{ACM-GlobalNorm} resolves this distortion by normalizing each task's projected Hessian contribution by its \emph{global trace} across all layers, preserving the relative sensitivity structure of the network layers. Although ACM-GlobalNorm yields a slightly lower overall average accuracy (\textbf{57.76\%}) compared to ACM-Norm (\textbf{58.89\%}) due to its global scaling constraint, it successfully avoids layer-wise trace-normalization distortions and outperforms ACM-Norm on individual tasks where global relative depth-wise sensitivity preservation is key (such as FashionMNIST at 70.02\% vs 69.24\% and CIFAR-10 at 77.05\% vs 76.27\%), thereby maintaining a physically consistent depth-wise sensitivity profile.
    \item \textbf{Bridging the Local-Global Gap and Sequential Coordination:} To practically bridge the local-global optimization gap, our Contracted ACM-GlobalNorm (using an unsupervised calibration-selected contraction factor of $\alpha=0.9$) achieves a Joint Average accuracy of \textbf{57.79\%}, slightly outperforming uncontracted ACM-GlobalNorm (\textbf{57.76\%}). This confirms that applying a slight global contraction is a highly viable zero-order regularization technique to pull merged parameters back into the valid local-quadratic surrogate regime. On the other hand, our Gauss-Seidel Coordinated ACM-GlobalNorm achieves a Joint Average accuracy of \textbf{36.65\%}, demonstrating that sequential coordination under coupled measurements is highly sensitive to cross-layer Hessian coupling block mismatches on deep networks.
\end{itemize}

\paragraph{Failure and Collapse of Test-Time Adaptation on Physical ViTs:}
While the state-of-the-art Test-Time Adaptation (TTA) methods (AdaMerging and PolyMerge) perform well in stylized simulations, they are severely compromised on physical Vision Transformer benchmarks, as quantitatively shown in Table \ref{tab:physical_results}. Under unsupervised entropy minimization on physical Vision Transformers without label supervision, the severe representational drift in the fused backbone results in local transductive overfitting to the small calibration sample. When we resolve the autograd graph disconnection of traditional TTA scripts, the optimizer aggressively overfits to the unsupervised target distribution, causing severe generalization collapse on the full test set. Specifically, AdaMerging achieves an average accuracy of \textbf{55.42\%}, PolyMerge collapses to \textbf{38.96\%}, and RegCalMerge (despite SNEW scale-weighting and ESR spatial regularization) is limited to \textbf{54.27\%}. This empirical degradation quantitatively demonstrates the extreme fragility and overfitting nature of test-time optimization heuristics on physical Vision Transformers, proving that analytical training-free frameworks like ACM are highly necessary.

\subsection{Layer-wise Analysis of Solved Coefficients}
To understand *how* ACM achieves this superior performance, we analyze the solved layer coefficients $\Lambda^{l, *}_{\text{Norm}}$ shown in Figure \ref{fig:lambdas} and detailed in Table \ref{tab:coefficients}.

\begin{table*}[ht]
\caption{Optimal layer coefficients $\Lambda^{l, *}_{\text{Norm}}$ solved by Scale-Normalized ACM (ACM-Norm) across the 14 parameter groups of the ViT-Tiny network.}
\label{tab:coefficients}
\vskip 0.15in
\begin{center}
\begin{small}
\begin{sc}
\begin{tabular}{lcccc}
\toprule
Layer Group $l$ & MNIST & FashionMNIST & CIFAR-10 & SVHN \\
\midrule
Layers 0 to 9 & 0.000 & 0.000 & 0.000 & 0.000 \\
Layer 10 & 0.147 & 0.524 & 0.597 & 0.269 \\
Layer 11 & 0.257 & 0.271 & 0.611 & 0.509 \\
Layer 12 & 0.393 & 0.483 & 0.616 & 0.454 \\
Layer 13 & 91.516 & -40.076 & 10.974 & -95.691 \\
\bottomrule
\end{tabular}
\end{sc}
\end{small}
\end{center}
\vskip -0.1in
\end{table*}

This analysis yields two extremely profound physical and mathematical insights:
\begin{enumerate}
    \item \textbf{Automatic Detection of Untrained Layers:} For Layers 0 to 9, which correspond to frozen and untrained layers in the experts, ACM automatically solves the coefficients as exactly $0.000$. This perfectly matches the physical reality: because the task vectors for these layers are of zero magnitude, their contribution to the loss landscape curvature is zero. While traditional test-time adaptation methods struggle to keep these coefficients stable (often introducing noise through gradients of unlabeled batches), ACM's analytical formulation is perfectly robust.
    \item \textbf{Interference Cancellation via Negative Coefficients:} In Layer 13 (the final LayerNorm), SVHN's coefficient is solved as $-95.691$, while FashionMNIST is solved as $-40.076$, MNIST is boosted to $91.516$, and CIFAR-10 is solved as $10.974$. This is a fascinating result: a negative coefficient represents an \emph{active cancellation mechanism}. Because task vectors are not orthogonal, the update directions of different tasks can interfere destructively. Under ACM, the off-diagonal terms of the projected Hessian $A^l_{ij}$ capture these cross-task directional alignments. In the final LayerNorm, the task vectors of FashionMNIST and other classification tasks exhibit a strong positive alignment. To prevent this destructive interference from shifting the final representations out of the valid classification regime, the solved analytical system automatically introduces a strong negative scaling for SVHN and FashionMNIST. This subtracts the interfering components and effectively "orthogonalizes" the update pathways at deployment. This is mathematically impossible to achieve in diagonal-approximation frameworks like Fisher Merging, highlighting the critical importance of full non-diagonal subspace curvature modeling.
\end{enumerate}

\begin{figure}[ht]
\vskip 0.2in
\begin{center}
\centerline{\includegraphics[width=\columnwidth]{acm_lambdas.png}}
\caption{Solved optimal layer-wise merging coefficients $\Lambda^{l, *}_{\text{Norm}}$ for each task using ACM-Norm on physical ViT-Tiny validation. Untrained frozen layers (0 to 9) are automatically set to 0.000, while active layers exhibit non-trivial scaling and negative coefficients that suppress cross-task representation interference.}
\label{fig:lambdas}
\end{center}
\vskip -0.2in
\end{figure}

\subsection{Discussion of Limitations, Hyperparameters, and Open Challenges}
To ensure absolute scientific transparency and provide actionable insights for future theoretical exploration, we dedicate this subsection to analyzing the structural limitations, hyperparameter behaviors, and mathematical trade-offs observed in our physical evaluation.

\paragraph{Practical Hyperparameter Selection at Deployment:}
A major concern in unsupervised model merging is the selection of the Ridge regularization scale $\gamma$ at deployment time without test-set leakage or label supervision. In practice, $\gamma$ can be selected in a completely unsupervised and training-free manner using only the tiny calibration dataset $\{\mathcal{D}_k\}$. Specifically, we propose a \emph{few-shot validation heuristic}: out of the 32 samples per task, 24 samples are allocated for projected Hessian/gradient estimation, and the remaining 8 samples are reserved as a local validation split. We sweep $\gamma$ over the candidate range $\{0.01, 0.05, 0.1, 0.2, 0.5, 1.0, 2.0\}$ and solve for the corresponding coefficients. We then evaluate the joint unsupervised prediction entropy (or Cross-Entropy Loss on the validation split) for each candidate $\gamma$, selecting the parameter that minimizes the validation loss. This unsupervised calibration-based check requires zero test-set access and takes less than 2 seconds, offering a robust deployment-ready pipeline.

\paragraph{The Local-Global Optimization Gap:}
A fundamental theoretical limitation of curvature-aware merging is the local-global approximation gap. ACM derives its analytical solution from local quadratic Taylor expansions taken around the individual task expert's minima $W_k$. However, the merged parameters $W(\Lambda)$ lie at a significant distance from any individual expert checkpoint (e.g., $W(\Lambda) - W_k \approx -0.7 v_k + \sum_{j \neq k} 0.3 v_j$). On highly non-convex, converged physical neural manifolds (such as our 10-epoch fine-tuned ViT backbones), the local quadratic Taylor approximation begins to break down over these large step sizes. This explains why on some settings, unguided baselines like standard Task Arithmetic (tuned to scale 0.4 at 60.72\%) can perform comparably to curvature-aware methods on highly non-convex, converged physical neural manifolds. In our current evaluation, our Vanilla ACM achieves 60.89\%, successfully outperforming Task Arithmetic (60.72\%), while ACM-GlobalNorm achieves 57.76\%. While Task Arithmetic lacks mathematical optimal guarantees, its uniform interpolation acts as a strong global regularizer that is robust to local landscape non-convexities. ACM, on the other hand, minimizes the local quadratic surrogate extremely precisely, but is slightly sub-optimal on the true global non-convex loss manifold under GlobalNorm traces.

In Appendix B.4, we derive a formal mathematical bound on this Taylor remainder (the local-global gap), proving that the approximation error is bounded by:
\begin{equation}
\begin{split}
    &\left| \mathcal{L}_k(W(\Lambda)) - \mathcal{L}_k^{\text{surr}}(W(\Lambda)) \right| \\
    &\le \frac{L_{\text{grad3}}}{6} \left( 1 + \|\Lambda\|_1 \right)^3 V_{\max}^3
\end{split}
\label{eq:main_local_global_bound}
\end{equation}
where $L_{\text{grad3}}$ is the Lipschitz constant of the loss third-order derivative, $V_{\max} = \max_{j} \left\| v_j \right\|_2$ is the maximum norm of the task vectors, and $\|\Lambda\|_1$ is the L1-norm of the merging coefficients. Equation \ref{eq:main_local_global_bound} shows that the local approximation error scales \emph{cubically} ($O(V_{\max}^3)$) with the magnitude of the task vectors. This mathematically formalizes why standard Task Arithmetic (which acts as a global uniform regularizer) can outperform local curvature-aware methods on highly fine-tuned physical models. This local-global optimization gap represents a landmark insight, suggesting that future research must incorporate higher-order tensor expansions (e.g., third-order derivatives) or zero-order validation corrections to bridge this gap.

\paragraph{Ill-Conditioning Risks and LayerNorm Distortion:}
Our layer-wise analysis revealed extremely large scaling and negative coefficients at Layer 13 (the final LayerNorm, e.g., SVHN solved to $-95.691$ and FashionMNIST to $-40.076$). While mathematically interpreted as an active cross-task representation cancellation mechanism, we must acknowledge that this is also an artifact of numerical ill-conditioning. Layer 13 represents the global LayerNorm, which has extremely few parameters (only 384 parameters in ViT-Tiny). Because of the low parameter count, the projected Hessian matrix $A^{13}$ has a very small trace and exhibits extreme ill-conditioning, with a condition number exceeding $10^4$. Solving the unregularized system results in massive coefficient blowups that degrade performance. The Ridge regularization parameter $\gamma$ acts as a crucial stabilizer by bounding the minimum eigenvalues of the regularized Hessian $A^l + \gamma I$, keeping the solved coefficients stable. 

To address this ill-conditioning without excessive scaling, we present a formal Lasso (L1) regularized proximal update scheme in Appendix B.6, which minimizes:
\begin{equation}
    \min_{\Lambda^l} \frac{1}{2} (\Lambda^l)^T A^l \Lambda^l - (\Lambda^l)^T (b^l - d^l) + \mu \left\| \Lambda^l \right\|_1
\label{eq:main_l1_obj}
\end{equation}
where $\mu > 0$ is the Lasso penalty strength. This objective can be solved efficiently using the proximal gradient method (Iterative Soft-Thresholding Algorithm, ISTA):
\begin{equation}
    \Lambda^{l,(s+1)} = \mathcal{S}_{\mu \eta} \left[ \Lambda^{l,(s)} - \eta \left( A^l \Lambda^{l,(s)} - (b^l - d^l) \right) \right]
\label{eq:main_ista}
\end{equation}
where $\mathcal{S}_{\tau}(x)_i = \text{sign}(x_i) \max\left( |x_i| - \tau, 0 \right)$ is the soft-thresholding operator and $\eta < 1 / \lambda_{\max}(A^l)$ is the learning rate. L1 regularization forces non-essential task vectors' coefficients to be exactly zero on highly ill-conditioned or low-parameter layers, preventing active cancellation blowups and acting as an automatic layer-wise expert task selector.

\paragraph{The Block-Jacobi Coupling Mismatch:}
Another elegant and subtle theoretical mismatch lies in the interaction between our decoupled Layer Block-Diagonal Assumption (Assumption 3.1) and our finite-difference calibration scheme. To decouple the optimization across layers, we assume that cross-layer second-order derivatives are negligible. However, during calibration, we perturb the weights of the \emph{entire} model simultaneously across all layers to estimate the perturbed gradient. Consequently, the measured gradient subtraction $g_{\text{sub}}^l = g_{k,j}^l - g_{k,0}^l$ at layer $l$ inherently captures the influence of perturbations in all other layers $l' \neq l$, effectively polling the cross-layer Hessian elements:
\begin{equation}
    g_{k,j}^l \approx g_{k,0}^l + \epsilon H_k^l v_j^l + \epsilon \sum_{l' \neq l} \left( \nabla^2_{W^l, W^{l'}} \mathcal{L}_k(W_k) \right) v_j^{l'}
\end{equation}
Despite obtaining these "coupling-polluted" measurements, ACM solves for the optimal coefficients layer-by-layer completely independently. This mathematical operation acts exactly like a single-step \emph{Block-Jacobi} solver applied to a coupled multi-layer system. While highly efficient and training-free, solving independently under coupled measurements introduces a projection error, especially on deep neural networks where early-layer scaling shifts late-layer representations. 

In Appendix B.5, we present a formal block Gauss-Seidel coordinate descent formulation to sequentially solve for coefficients while holding other layers fixed. Under this scheme, for each layer $l$, the coefficient vector is iteratively updated as:
\begin{equation}
\begin{split}
    \left( A^{l,l} + \gamma I \right) \Lambda^{l,(t)} &= b^l - d^l - \sum_{l' < l} A^{l, l'} \Lambda^{l',(t)} \\
    &- \sum_{l' > l} A^{l, l'} \Lambda^{l',(t-1)}
\end{split}
\label{eq:main_block_gs}
\end{equation}
where $A^{l, l'}$ represents the cross-layer projected Hessian block coupling layers $l$ and $l'$.

Our physical evaluation in Table \ref{tab:physical_results} shows that this sequential Gauss-Seidel coordination drops to \textbf{36.65\%} Joint Average accuracy, which serves as a highly valuable diagnostic insight into the step-by-step mechanics of this collapse:
\begin{enumerate}
    \item \textbf{Representational Drift Cascade:} When sequentially updating coefficients layer-by-layer starting from the early layers, small approximation errors in early blocks (e.g., layers 10--11) accumulate. Because these layers control the input projections and early-feature transformations, even minor deviations in their coefficients alter the downstream activation features passed to subsequent layers.
    \item \textbf{Hessian Reference-Point Collapse:} In standard Taylor expansions, the Hessian of a downstream layer (e.g., Layer 12) is evaluated strictly at the original, unperturbed activation manifold of the task experts. Once the early layers are modified sequentially, the activations arriving at downstream blocks reside in a completely different, shifted region of the representation space. Consequently, the local quadratic surrogate computed around the \emph{original} unperturbed expert checkpoints becomes invalid, causing the downstream optimization to collapse.
    \item \textbf{Out-of-Distribution LayerNorm Blowup:} LayerNorm layers scale inputs by the inverse of their standard deviation. When representation distributions drift due to sequential early-layer updates, the variance of downstream activations fluctuates wildly. This triggers severe numerical ill-conditioning in the downstream LayerNorm parameters, amplifying small scaling errors into catastrophic representational divergence.
    \item \textbf{Non-Contraction Operator Divergence:} Because the cross-layer Hessian coupling blocks $A^{l, l'}$ are approximated via finite-difference measurements, the block Gauss-Seidel update operator does not form a contraction mapping on non-convex physical manifolds. Instead of converging to a stable joint minimum, sequential updates propagate errors exponentially across layers, causing the model to get trapped in extremely sharp, sub-optimal local minima.
\end{enumerate}
In contrast, the simultaneous Block-Jacobi analytical solution (Vanilla ACM) acts as a global, single-step projection that remains surprisingly regularized and stable. This diagnostic indicates that future multi-layer coordination must be solved dynamically or using soft relaxation. Recognizing this coupling mismatch highlights an exciting theoretical direction: future work can explore iterative block-Jacobi relaxation schemes or coordinate descent over layers to fully resolve the multi-layer Hessian coupling.

\paragraph{Scaling Analysis and Architectural Diversity:}
To address the limited architectural and domain diversity in our physical validation, we provide a formal scaling analysis of ACM's numerical behavior under larger model backbones (e.g., ViT-Base, ViT-Large, or RoBERTa/LLaMA for NLP). When scaling the backbone architecture, two primary geometric transformations occur on the neural manifold: (1) The hidden/embedding dimension $D_{\text{embed}}$ scales up (e.g., from 192 in ViT-Tiny to 768 in ViT-Base and 1024 in ViT-Large). Consequently, the parameter count of previously low-parameter bottleneck layers (like LayerNorm layers) increases linearly with the dimension (e.g., from 384 parameters to 1536 and 2048 parameters). This substantially increases the trace of their projected Hessians and reduces the corresponding condition number by several orders of magnitude, effectively mitigating LayerNorm ill-conditioning and eliminating coefficient blowups without requiring aggressive regularization. (2) As model capacity increases, the task-specific update directions (task vectors) become highly orthogonal within the massive parameter space. This intrinsic high-dimensional sparsity reduces destructive cross-task representational interference, which in turn reduces the magnitude of the off-diagonal coupling terms $A^{l, l'}$. This means that on larger models, the block-Jacobi projection error naturally shrinks, making local curvature-based analytical merging even more accurate and stable. This scaling behavior indicates that ACM's performance is expected to improve as models scale, establishing a promising theoretical trajectory for future large-scale weight consolidation.

In summary, ACM is not only mathematically elegant but also exhibits deep physical compliance with the underlying geometry of neural networks. We hope that this rigorous characterization of its limitations stimulates deeper theoretical research in the weight consolidation community.
